% =====================================================================
% Section 3: System Architecture
% Overleaf 可直接粘贴。要求：LNCS (llncs.cls)。
% 粘贴位置：Section 2 (Data Resource) 之后、Demonstration 之前。
% 注意：
%   1) Fig.1 文件：out/fig1_firstcare_architecture.svg 需先转为 PDF
%      （Inkscape / rsvg-convert / 在线工具），放入工程后改下方文件名；
%   2) \ref{sec:data} 对应 Section 2 的 \label{sec:data}；
%   3) 数学记法与 Fig.1 严格一致：$p_{\mathrm{men}}(f)$、
%      $p_{\mathrm{sup}}(f)$、$\tau_f = 0.7$。
% =====================================================================

\section{System Architecture}
\label{sec:architecture}

\begin{figure}[t]
\centering
\includegraphics[width=\textwidth]{fig1_firstcare_architecture_3x.png}%
% 已用 out/fig1_firstcare_architecture_3x.png（4080x1770，方案 A 大字号版，
% 印刷 12.2x5.29cm，849 DPI）。若后续转出矢量 PDF 可替换以获得更佳印刷质量。
\caption{Overview of FirstCare. A colloquial patient query is matched against precomputed atomic propositions through two parallel questions; each proposition is classified as \textit{supported}, \textit{gap}, or \textit{absent} with threshold $\tau_f=0.7$; \textit{absent} propositions are removed by a hard gate before grounded triage generation, and \textit{gap}/\textit{absent} facts are returned as follow-up chips for multi-turn clarification. HPI: history of present illness; Dx/DDx: diagnosis and differential diagnosis.}
\label{fig:architecture}
\end{figure}

As shown in Fig.~\ref{fig:architecture}, FirstCare consists of three stages: hybrid candidate retrieval, dual-question fact verification, and grounded triage generation, closed by a multi-turn clarification loop in which unverified propositions are surfaced as follow-up chips; user-supplied statements are accumulated with the original query and the full pipeline is re-run.

\subsection{Hybrid Candidate Retrieval}

Each structured case is indexed as a standardized document view (title, patient general information including the history of present illness, and diagnosis with differential diagnosis) using character-level 1/2/3-gram BM25 and TF-IDF cosine similarity with equal weights, yielding the top-50 candidates. A discriminative verification model, Jev,\footnote{Jev is a discriminative System One model that takes a state and typed questions and returns type-safe, calibrated probability estimates without generating text; see \url{https://docs.typesafe.ai/concepts/system-one}.} then reranks the candidates pointwise and returns the top-$k$ candidates ($k=10$). Under a stratified 5-fold cross-validation protocol, the hybrid retriever achieves a mean Recall@50 of 0.941. We also report a negative result: LoRA fine-tuning of a pretrained sentence-embedding model on this small corpus causes catastrophic forgetting (Recall@50 of 0.237 on fold~0, below both the zero-shot backbone at 0.467 and the hybrid retriever at 0.948); dense retrieval is therefore kept as an ablation rather than the main line.

\subsection{Dual-Question Fact Verification}

The patient input first undergoes deterministic colloquial normalization: a lexicon of nine categories maps metaphors and colloquial phrasings to standard clinical terms, and a contrast line is appended to the verification input without altering the patient's original wording. Rather than verifying whole sentences, the system extracts a core element phrase from the mention keywords of each precomputed atomic proposition (Sect.~\ref{sec:data}) and formulates element-level questions, to which Jev returns two calibrated probabilities in parallel: the mention probability $p_{\mathrm{men}}(f)$, whether the patient has stated element~$f$, and the support probability $p_{\mathrm{sup}}(f)$, whether $f$ is supported by the patient's statements. Each proposition is then assigned one of three states with threshold $\tau_f=0.7$:
\begin{equation}
\mathrm{state}(f)=
\begin{cases}
\textit{supported}, & p_{\mathrm{sup}}(f)\ge\tau_f,\\[2pt]
\textit{gap}, & p_{\mathrm{sup}}(f)<\tau_f \,\wedge\, p_{\mathrm{men}}(f)\ge\tau_f,\\[2pt]
\textit{absent}, & p_{\mathrm{men}}(f)<\tau_f.
\end{cases}
\label{eq:threestate}
\end{equation}
Propositions in the \textit{gap} and \textit{absent} states are exported as follow-up chips; after the user answers, the accumulated statements re-run the entire pipeline, realizing the interactive clarification loop.

\subsection{Grounded Triage Generation}

Before generation, a hard gate removes all \textit{absent} propositions, guaranteeing that no unstated symptom appears in the output. If any red-flag proposition is \textit{supported}, a red-flag escalation rule enforces an immediate emergency-care recommendation. The generator produces a five-step structured output grounded in the top-$k$ candidate cases: triage conclusion, misconception correction, fact mapping, boundary conditions, and next-step advice. It operates in two modes: a template-based deterministic mode (default), which is fully reproducible, and an LLM mode whose output is adopted only after passing red-flag and structure validation, with any failure reverting to the deterministic template. Four additional safeguards are applied: a specificity check on numeric and temporal qualifiers; a split of \textit{supported} propositions into two columns, \textit{patient-confirmed} versus \textit{case-derived features pending verification}, with only the former counted as patient-confirmed facts; a recall-oriented red-flag rescue mechanism that prefers escalation over omission; and a two-layer post-hoc fact check that detects case-information leakage and fabricated facts, rewriting flagged phrases as case-derived features pending verification.
